\documentclass[11pt]{article}
\usepackage[margin=0.7in]{geometry}
\usepackage{amsmath,amssymb}
\usepackage{enumitem}
\usepackage{tikz}
\usepackage{pgfplots}
\pgfplotsset{compat=1.18}
\usepackage{fancyhdr}
\usepackage{parskip}
\setlist[enumerate]{leftmargin=*, itemsep=6pt}
\pagestyle{fancy}
\fancyhf{}
\fancyfoot[L]{\small\textit{en}Vision\textsuperscript{\textregistered} Algebra 2}
\fancyfoot[C]{\small 2-4 Lesson Quiz, short form}
\fancyfoot[R]{\small\thepage}
\renewcommand{\headrulewidth}{0pt}
\newcommand{\choice}[1]{\tikz[baseline=-0.6ex]\node[draw, circle, inner sep=0.6pt, minimum size=3.3mm, font=\scriptsize\bfseries] {#1};\;}
\newcommand{\checkbox}{$\square$\;}
\pgfplotsset{quizgrid/.style={
  axis lines=middle, grid=both, grid style={gray!45},
  ticklabel style={font=\scriptsize, fill=white, inner sep=1pt},
  xlabel={$x$}, ylabel={$y$}, axis line style={thick}, clip=true}}

\begin{document}

\noindent\textbf{Name}\; \rule{6cm}{0.4pt} \hfill Period \rule{1.2cm}{0.4pt} \quad Date \rule{2.2cm}{0.4pt}

\vspace{4pt}\noindent\rule{\textwidth}{0.6pt}

\vspace{2pt}\noindent{\LARGE\textbf{2-4 \; Lesson Quiz: Complex Numbers and Operations}} \hfill {\small short form}

\vspace{4pt}\noindent\rule{\textwidth}{0.6pt}

\noindent{\small\textit{Questions 2, 3 and 4 of the Savvas 2-4 Lesson Quiz, numbered as on the Savvas form. Show your work.}}

\vspace{6pt}
\begin{minipage}[t]{0.50\textwidth}
\textbf{2.} What is the difference of $(5+3i)$ and $(2+9i)$?

\vspace{4pt}
\begin{itemize}[leftmargin=1.2em, itemsep=5pt]
\item[] \choice{A} $7+12i$
\item[] \choice{B} $3+6i$
\item[] \choice{C} $3-6i$
\item[] \choice{D} $7+6i$
\end{itemize}
\end{minipage}\hfill
\begin{minipage}[t]{0.46\textwidth}
\textit{Work:}

\vspace{30pt}\rule{\linewidth}{0.4pt}

\vspace{18pt}\rule{\linewidth}{0.4pt}
\end{minipage}

\vspace{10pt}\noindent\rule{\textwidth}{0.4pt}

\vspace{6pt}
\begin{minipage}[t]{0.50\textwidth}
\textbf{3.} Write the product $(2+7i)(2-7i)$ in the form $a+bi$.

\vspace{4pt}
\begin{itemize}[leftmargin=1.2em, itemsep=5pt]
\item[] \choice{A} $4-49i$
\item[] \choice{B} $53$
\item[] \choice{C} $4-14i^2$
\item[] \choice{D} $53+14i$
\end{itemize}
\end{minipage}\hfill
\begin{minipage}[t]{0.46\textwidth}
\textit{Work:}

\vspace{30pt}\rule{\linewidth}{0.4pt}

\vspace{18pt}\rule{\linewidth}{0.4pt}

\vspace{18pt}\rule{\linewidth}{0.4pt}
\end{minipage}

\vspace{10pt}\noindent\rule{\textwidth}{0.4pt}

\vspace{6pt}
\textbf{4.} Write the expression $3i(2-5i)(4+i)$ in the form $a+bi$.

\vspace{6pt}
\textit{Work:}

\vspace{26pt}\rule{\linewidth}{0.4pt}

\vspace{20pt}\rule{\linewidth}{0.4pt}

\vspace{20pt}\rule{\linewidth}{0.4pt}

\vspace{20pt}\rule{\linewidth}{0.4pt}

\vspace{14pt}
\hfill $\rule{1.6cm}{0.4pt}\;+\;\rule{1.6cm}{0.4pt}\,i$

\vfill
\noindent\textbf{Scoring:} Question 3 and Question 4 are 10 points each. Question 2 is a bonus of 5 points. Total out of 20.

\newpage
\noindent{\LARGE\textbf{2-4 Lesson Quiz, short form: answer key}}

\vspace{4pt}\noindent\rule{\textwidth}{0.6pt}

\begin{enumerate}
\item[2.] \textbf{C, $3-6i$.} \mbox{$(5+3i)-(2+9i)=(5+3i)+(-2-9i)=(5-2)+(3i-9i)=3-6i$}. A adds the two numbers; B subtracts the imaginary parts in the wrong order ($9i-3i$); D adds the real parts and makes the same imaginary-part error. Savvas quiz item 2, DOK 1, Skills Review NS11 (taught by Example 2, cut from the fast packet; Practice b rehearses it). \textbf{Bonus 5:} 5 for C with the subtraction shown; 3 for C with no work; 0 otherwise.
\item[3.] \textbf{B, $53$.} \mbox{$(2+7i)(2-7i)=4-14i+14i-49i^2=4-49(-1)=53$}. The factors are complex conjugates, so the product is $2^2+7^2=53$, a real number. A writes $-49i^2$ as $-49i$; C and D mishandle the middle terms or $i^2$. Savvas quiz item 3, DOK 1, Skills Review NS11 (Example 3B, Try It 3b). \textbf{10 points:} 7 for B, 3 for work that shows $i^2=-1$ used.
\item[4.] \textbf{$54+39i$} (blanks: 54; 39).\\
$(2-5i)(4+i)=8+2i-20i-5i^2=8-18i+5=13-18i$, then $3i(13-18i)=39i-54i^2=54+39i$.\\
Other order, same answer: $3i(2-5i)=15+6i$, then $(15+6i)(4+i)=60+15i+24i+6i^2=54+39i$.\\
Savvas quiz item 4, DOK 2, Skills Review NS11 (Example 3A and 3B; Practice a is the parallel). \textbf{10 points:} 4 for the first product, 4 for the second product with $i^2=-1$, 2 for the answer written as $a+bi$.
\end{enumerate}

\textbf{Items left off this form} (available for a retake or Form B): Savvas item 1, solve $x^2=-64$ over the complex numbers, select all (answer A and D, $8i$ and $-8i$; DOK 1; Example 1, cut); Savvas item 5, match each product with its factors (answers: $18x^2+2$ is $2$, $(3x-i)$, $(3x+i)$; $3x^2+12$ is $3$, $(x-2i)$, $(x+2i)$; $36x^2+25$ is $(6x-5i)(6x+5i)$; $9x^2+1$ is $(3x-i)(3x+i)$; DOK 2; Example 5, cut). Both use the Skills Review code NS11. The TE prints no standards codes; the curriculum map's Topic 2 list gives HSN.CN.A.1 and HSN.CN.A.2 for this lesson's operations.

\textbf{IXL replacement:} SmartScore 70 or above on each starred skill, VZ8, 7U5 and JVF, replaces this quiz with 20/20.

\end{document}
